\begin{figure}[ht]
  \centering
  \begin{tikzpicture}
    \begin{groupplot}[group style={group size=3 by 1, horizontal sep=1.3cm},
        tesi, width=.32\textwidth, height=5.4cm, xmin=-.3, xmax=7.3,
        xtick={0,...,7}, xticklabels={0,5,10,20,40,80,160,320},
        x tick label style={font=\tiny, rotate=90, anchor=east}, xlabel={costante}]
      \nextgroupplot[title={\footnotesize SciFact, nDCG@10}]
        \addplot[color=black, mark=*, mark size=1.5pt] table[x=posizione, y=scifact] {dati/fusione.dat};
      \nextgroupplot[title={\footnotesize NFCorpus, nDCG@10}]
        \addplot[color=black, mark=*, mark size=1.5pt] table[x=posizione, y=nfcorpus] {dati/fusione.dat};
      \nextgroupplot[title={\footnotesize known-item, MRR@10}]
        \addplot[color=black, mark=*, mark size=1.5pt] table[x=posizione, y=known] {dati/fusione.dat};
    \end{groupplot}
  \end{tikzpicture}
  \caption[La costante della somma pesata sulle collezioni di calibrazione]{La costante della somma pesata al variare del suo valore, sulle
  tre collezioni di calibrazione (SciFact e known-item, sezione di addestramento; NFCorpus, sviluppo). Il valore di
  partenza di Koskidex è 20: vicino all'ottimo sulle known-item, e otto volte più
  basso dell'ottimo sulle query in lingua naturale. Dati: \texttt{latex/dati/fusione.dat}.}
  \label{fig:fusione}
\end{figure}
